\chapter{Evaluation der Hypothesen}
\label{ch:evaluation}

In diesem Kapitel werden die in Abschnitt \ref{sec:hypothesen} aufgestellten
Hypothesen anhand der implementierten Prototypen evaluiert. Die Messung erfolgt
nach den in Abschnitt \ref{sec:evaluationsrahmen} definierten Kriterien und
orientiert sich an der evaluativen Case Study nach \citep{Kitchenham1995}. Die
Ergebnisse werden in Tabellen dokumentiert, um einen direkten Vergleich
zwischen \ac{AWS} Cognito mit Spring Boot und \ac{SAP IAS} mit \ac{SAP CAP} zu
ermöglichen.


\section{Übersicht der Messergebnisse}

Tabelle \ref{tab:hypothesen_uebersicht} zeigt eine Zusammenfassung aller
Hypothesen mit den gemessenen Metriken und dem Evaluationsergebnis. Die
Hypothesen sind in zwei Kategorien unterteilt: Framework-Integration ($H_1$ bis
$H_5$) und direkter IdP-Vergleich ($H_6$ bis $H_9$).

\begin{table}[h]
    \centering
    \small
    \begin{tabular}{|l|p{4cm}|c|c|c|}
        \hline
        \textbf{H}  & \textbf{Beschreibung} & \textbf{AWS/Spring}      & \textbf{SAP CAP/IAS}    & \textbf{Ergebnis} \\
        \hline
        \multicolumn{5}{|c|}{\textbf{Framework-Integration}}                                                                           \\
        \hline
        $H_1$ & Autorisierung (CC + Dateien)  & CC=26, 5 Dateien          & CC=0, 2 Dateien & SAP: CC=0        \\
        \hline
        $H_2$ & Wartbarkeit (Dateien bei Rollenänderung)  & 3-4 Dateien, 5 Methoden  & 1 Datei, 0 Methoden     & SAP: -75\%        \\
        \hline
        $H_{3a}$  & Token-Validierung (LOC + Komplexität) & 150 LOC, CC=12           & 10 LOC, CC=0            & SAP: -93\%        \\
        \hline
        $H_{3b}$  & Unternehmensintegration (Schritte)  & 6 Schritte                     & 8 Schritte                    & AWS: -25\%              \\
        \hline
        $H_4$ & Vendor-Kopplung (Dateien) & 5 Dateien                & 2 Dateien               & SAP: -60\%        \\
        \hline
        $H_5$ & Token-Validierung (Komplexität) & CC=12, manuell           & CC=0, automatisch       & SAP: CC=0         \\
        \hline
        \multicolumn{5}{|c|}{\textbf{Direkter IdP-Vergleich}}                                                                          \\
        \hline
        $H_6$ & Benutzerregistrierung (Klicks)  & 5 Klicks                     & 4 Klicks                    & Vergleichbar              \\
        \hline
        $H_7$ & Recovery-Mechanismen (Sicherheit) & Code-basiert             & Link-basiert            & SAP: sicherer              \\
        \hline
        $H_8$ & Integrierte Sicherheitsfeatures & 2 Standard, 5 konfig.                     & 6 Standard, 4 konfig.                    & SAP: +200\%              \\
        \hline
        $H_9$ & Benutzerlöschung (Klicks) & 6 Klicks                     & 5 Klicks                    & Vergleichbar              \\
        \hline
    \end{tabular}
    \caption{Übersicht der Hypothesen-Evaluation}
    \label{tab:hypothesen_uebersicht}
\end{table}


\section{H1: Sicherheit der Autorisierungslogik}

Hypothese $H_1$ postuliert, dass der deklarative Sicherheitsansatz von \ac{SAP
CAP} die Anzahl notwendiger Autorisierungs-Code-Zeilen gegenüber dem
imperativen Spring Security Ansatz reduziert.

\subsection{Messmethodik}

Die Messung erfolgt durch statische Code-Analyse der zyklomatischen Komplexität (CC) nach \citep{McCabe1976}. Die CC quantifiziert die Anzahl unabhängiger Pfade durch den Code: CC = E - N + 2P (E=Kanten, N=Knoten, P=Komponenten). In der Praxis entspricht dies der Anzahl der Entscheidungspunkte (if, for, while, catch, \&\&, ||) plus 1.
\begin{itemize}
    \item \textbf{Spring Boot:} CC der Methoden in \texttt{JwtAuthenticationFilter}, \texttt{SecurityConfig} und \texttt{CognitoService.validateToken()}
    \item \textbf{SAP CAP:} CC der \texttt{@restrict}-Annotationen im \ac{CDS}-Modell (deklarativ = CC=0)
\end{itemize}

\subsection{Messergebnisse}

\begin{table}[h]
    \centering
    \begin{tabular}{|l|c|c|}
        \hline
        \textbf{Metrik}          & \textbf{AWS Cognito + Spring Boot} & \textbf{SAP CAP + SAP IAS}  \\
        \hline
        \multicolumn{3}{|c|}{\textbf{Zyklomatische Komplexität (CC)}} \\
        \hline
        JwtAF.doFilterInternal                     & 3              & --  \\
        \hline
        JwtAF.extractJwtFromRequest                & 3                    & --  \\
        \hline
        JwtAF.extractAuthorities                   & 6                & --  \\
        \hline
        CognitoService.validateToken     & 13             & --  \\
        \hline
        SC.securityFilterChain                 & 1              & --  \\
        \hline
        data-model.cds @restrict     & --             & 0 \\
        \hline
        \textbf{Gesamt CC} & \textbf{26}              & \textbf{0}  \\
        \hline
        \multicolumn{3}{|c|}{\textbf{Weitere Metriken}} \\
        \hline
        @PreAuthorize Annotationen & 4              & 0 \\
        \hline
        @restrict Annotationen & 0              & 4 \\
        \hline
        Anzahl Security-Dateien  & 5              & 2 \\
        \hline
    \end{tabular}
    \footnotesize\\
    \textbf{Legende:} SC = SecurityConfig, JwtAF = JwtAuthenticationFilter
    \caption{Messergebnisse für H1: Autorisierungslogik -- Zyklomatische Komplexität}
    \label{tab:h1_ergebnisse}
\end{table}

\subsection{Interpretation}

Die Ergebnisse zeigen einen fundamentalen Unterschied zwischen den beiden Ansätzen.
Bei Spring Boot beträgt die zyklomatische Komplexität der Autorisierungslogik CC=26, verteilt über 5 Dateien mit 12 Entscheidungspunkten allein in der Token-Validierung. Bei \ac{SAP CAP} beträgt die zyklomatische Komplexität CC=0, da die \texttt{@restrict}-Annotationen rein deklarativ sind und keine programmatische Verzweigungslogik enthalten.
Die hohe zyklomatische Komplexität bei Spring Boot erhöht das Risiko für Implementierungsfehler gemäß \ac{OWASP} CNAS-3 (\textit{Improper Authentication \& Authorization}). Jeder zusätzliche Entscheidungspunkt ist eine potenzielle Fehlerquelle für Sicherheitslücken. $H_1$ wird bestätigt: Der deklarative Ansatz von \ac{SAP CAP} eliminiert die zyklomatische Komplexität der Autorisierungslogik vollständig (CC=26 $\rightarrow$ CC=0).


\section{H2: Wartbarkeit bei Rollenänderungen}

Hypothese $H_2$ untersucht, ob Änderungen an der Autorisierungslogik im \ac{SAP
CAP} Ökosystem weniger Quellcode-Änderungen erfordern.

\subsection{Testszenario}

Für den Test wurde eine neue Rolle \texttt{Observer} mit Leseberechtigung
hinzugefügt. Gemessen wurde:
\begin{itemize}
    \item Anzahl der geänderten Dateien
    \item Anzahl der geänderten Methoden
    \item Gesamte LOC-Änderungen
\end{itemize}

\subsection{Messergebnisse}

\begin{table}[h]
    \centering
    \begin{tabular}{|l|c|c|}
        \hline
        \textbf{Metrik} & \textbf{AWS Cognito + Spring Boot}  & \textbf{SAP CAP + SAP IAS}  \\
        \hline
        Geänderte Dateien & 3-4 & 1 \\
        \hline
        Geänderte Methoden  & 5 & 0 \\
        \hline
        LOC-Änderungen  & $\sim$18  & 3 \\
        \hline
    \end{tabular}
    \caption{Messergebnisse für H2: Wartbarkeit}
    \label{tab:h2_ergebnisse}
\end{table}

\subsection{Interpretation}

Bei \ac{SAP CAP} wurde die Observer-Rolle durch Hinzufügen von 3 Zeilen in der
\texttt{data-model.cds} implementiert. Keine Java-Datei musste geändert werden.
Bei Spring Boot müssen dagegen mehrere Controller-Methoden mit neuen
\texttt{@PreAuthorize}-Annotationen versehen werden. Dies bestätigt $H_2$: Der
deklarative Ansatz verbessert die Wartbarkeit bei Rollenänderungen signifikant
(-75\% Dateien, -100\% Methoden).


\section{H3a: Token-Validierung und JWT-Handling}

Hypothese $H_{3a}$ vergleicht den Code- und Konfigurationsaufwand für
Token-Validierung. Zusätzlich zu den Lines of Code wird die zyklomatische
Komplexität gemessen.

\subsection{Messergebnisse}

\begin{table}[h]
    \centering
    \begin{tabular}{|l|c|c|}
        \hline
        \textbf{Komponente} & \textbf{AWS Cognito + Spring Boot}  & \textbf{SAP CAP + SAP IAS}  \\
        \hline
        JwtAuthenticationFilter & 110 LOC, CC=5 & 0 \\
        \hline
        CognitoService.validateToken()  & 60 LOC, CC=7  & 0 \\
        \hline
        application.yaml (Security) & 0 & 10 LOC  \\
        \hline
        \textbf{Gesamt LOC} & $\sim$150 & $\sim$10  \\
        \hline
        \textbf{Zyklomatische Komplexität}  & CC=12 & CC=0  \\
        \hline
    \end{tabular}
    \caption{Messergebnisse für H3a: Token-Validierung}
    \label{tab:h3a_ergebnisse}
\end{table}

\subsection{Interpretation}

\ac{SAP CAP} übernimmt die Token-Validierung vollständig durch das Framework. Die einzige Konfiguration erfolgt in der \texttt{application.yaml} durch Aktivierung des XSUAA-Modus für die Cloud-Umgebung. Bei Spring Boot ist dagegen manuelle JWKS-Konfiguration, Signaturverifikation und Claims-Validierung erforderlich. Die zyklomatische Komplexität von CC=12 bei Spring Boot gegenüber CC=0 bei \ac{SAP CAP} zeigt das erhöhte Fehlerpotenzial des manuellen Ansatzes. $H_{3a}$ wird bestätigt: \ac{SAP CAP} reduziert die Token-Validierungs-Komplexität um 93\% (LOC) bzw. 100\% (Komplexität).


\section{H3b: Unternehmensintegration}

Hypothese $H_{3b}$ misst den Konfigurationsaufwand bei der Anbindung eines
externen \ac{IdP}.

\subsection{Messmethodik}

Die Messung erfolgt in der Cloud-Umgebung:
\begin{itemize}
    \item \textbf{SAP BTP:} Konfigurationsschritte für Service-Binding zu \ac{SAP IAS}
    \item \textbf{AWS:} Konfigurationsschritte für Cognito User Pool Federation
\end{itemize}

\subsection{Messergebnisse}

Als Ausgangspunkt wird eine beispielsweise eine Microsoft Entra ID als \ac{IdP} angenommen, welcher hinzugefügt werden soll.
An den Backends ist keine Code-Änderungen notwendig, nur Konfigurationsschritte in der Cloud, jedoch können URL-Änderungen bei \ac{AWS} Cognito erforderlich sein, dies ist abhängig von der Konfiguration.
\ac{AWS} Cognito kann direkt mit dem User Pool verbunden werden, während bei \ac{SAP IAS} ein Service-Binding in der SAP BTP erstellt werden muss.
Daher bietet \ac{AWS} Cognito mit 6 Konfigurationsschritten weniger als \ac{SAP IAS} mit 6 Konfigurationsschritten.
Abbildung \ref{fig:aws_extidp} zeigt die Konfiguration eines externen \ac{IdP} in \ac{AWS} Cognito. Dabei können auch unterschiedliche Standard-Protokolle wie \ac{SAML} 2.0 oder OpenID Connect verwendet werden.

\begin{figure}[h]
    \centering
    \includegraphics[width=0.95\textwidth]{figures/aws_extidp.png}
    \caption{Screenshot aus AWS Cognito zur Konfiguration eines externen IdP}
    \label{fig:aws_extidp}
\end{figure}


\section{H4: Vendor-Abhängigkeit und Framework-Kopplung}

Hypothese $H_4$ quantifiziert den Kopplungsgrad zwischen Applikationscode und
Framework.

\subsection{Messmethodik}

Gezählt werden Dateien mit Framework-spezifischen Imports:
\begin{itemize}
    \item \textbf{Spring Boot:} Dateien mit \texttt{org.springframework.security.*} Imports
    \item \textbf{SAP CAP:} Dateien mit \texttt{@sap/*} Imports
\end{itemize}

\subsection{Messergebnisse}

\begin{table}[h]
    \centering
    \begin{tabular}{|l|c|c|}
        \hline
        \textbf{Metrik} & \textbf{AWS Cognito + Spring Boot}  & \textbf{SAP CAP + SAP IAS}  \\
        \hline
        Dateien mit Security-Imports  & 5 & 2 \\
        \hline
        Gekoppelte Pakete & 3 (AWS SDK, Spring Security, Nimbus JWT)  & 2 (SAP CDS, Spring Security Core) \\
        \hline
        @PreAuthorize in Controllern  & 4 & 0 \\
        \hline
        @restrict in CDS  & 0 & 4 \\
        \hline
    \end{tabular}
    \caption{Messergebnisse für H4: Vendor-Kopplung}
    \label{tab:h4_ergebnisse}
\end{table}

\subsection{Interpretation}

Die detaillierte Analyse zeigt, dass Spring Boot 5 Java-Dateien mit Security- oder IdP-spezifischen Imports benötigt: \texttt{SecurityConfig}, \texttt{JwtAuthenticationFilter}, \texttt{CognitoService}, \texttt{OrderController} und \texttt{GlobalExceptionHandler}. Dabei ist besonders relevant, dass der \texttt{OrderController} als Business-Logik-Datei Spring Security Imports für \texttt{@PreAuthorize}-Annotationen benötigt.

\ac{SAP CAP} benötigt nur 2 Dateien mit Security-bezogenem Code: \texttt{OrderServiceHandler.java} (für \texttt{SecurityContextHolder} zur Username-Extraktion) und \texttt{data-model.cds} (für \texttt{@restrict}-Annotationen). Die Java-Datei enthält keine Autorisierungslogik, diese wird vollständig durch das \ac{CDS}-Modell durchgesetzt.

Die Reduktion von 60\% (5 $\rightarrow$ 2 Dateien) bestätigt $H_4$: \ac{SAP CAP} reduziert die Vendor-Kopplung signifikant. Der entscheidende Unterschied liegt darin, dass bei \ac{SAP CAP} die Business-Logik-Dateien keine Security-Annotationen benötigen.


\section{H5: Automatische Token-Validierung}

Hypothese $H_5$ untersucht, ob die Token-Validierung automatisch oder manuell
erfolgt.

\subsection{Messergebnisse}

\begin{table}[h]
    \centering
    \begin{tabular}{|l|c|c|}
        \hline
        \textbf{Aspekt}  & \textbf{AWS Cognito + Spring Boot}  & \textbf{SAP CAP + SAP IAS} \\
        \hline
        Validierung  & Manuell (JwtAuthenticationFilter) & Automatisch (Framework)  \\
        \hline
        JWKS-Konfiguration & Manuell (application.yaml)  & Automatisch (Service-Binding)  \\
        \hline
        Zyklomatische Komplexität  & CC=12 & CC=0 \\
        \hline
        LOC für Validierung  & $\sim$150 & 0  \\
        \hline
    \end{tabular}
    \caption{Messergebnisse für H5: Token-Validierung}
    \label{tab:h5_ergebnisse}
\end{table}

\subsection{Interpretation}

\ac{SAP CAP} validiert Tokens vollständig automatisch durch die \texttt{cds-feature-identity} Bibliothek. In der Cloud erfolgt die Konfiguration über Service-Binding, ohne manuelle JWKS-URI-Angabe. Spring Boot erfordert hingegen eine explizite Konfiguration und einen manuellen JWT-Filter mit einer zyklomatischen Komplexität von CC=12. Diese Komplexität erhöht das Fehlerpotenzial bei Sicherheits-kritischem Code. $H_5$ wird bestätigt: \ac{SAP CAP} bietet vollautomatische Token-Validierung mit CC=0.


\section{H6: Benutzerregistrierung}

Hypothese $H_6$ vergleicht den Aufwand für Benutzer-Onboarding durch Messung
der Klicks in beiden Admin-Konsolen. Dies ist ein direkter \ac{IdP}-Vergleich,
unabhängig von der Framework-Implementierung.

\subsection{Messmethodik}

Die Benutzererstellung wird in beiden Admin-Konsolen durchgeführt und die
Anzahl der Klicks/Schritte gezählt.

\subsection{Messergebnisse: Admin-Konsolen-Vergleich}

\begin{table}[h]
    \centering
    \begin{tabular}{|l|c|c|}
        \hline
        \textbf{Schritt}            & \textbf{AWS Cognito}          & \textbf{SAP IAS}          \\
        \hline
        1. Admin-Konsolen/UI öffnen & 1                            & 1                        \\
        \hline
        2. User Pool / Application auswählen            & 1     & 0     \\
        \hline
        3. Users-Bereich navigieren           & 1     & 1     \\
        \hline
        4. "Create User" klicken            & 1     & 1     \\
        \hline
        5. Formular ausfüllen           & 0     & 0     \\
        \hline
        6. Bestätigen           & 1     & 1     \\
        \hline
        \textbf{Gesamt Klicks}            & 5     & 4     \\
        \hline
    \end{tabular}
    \caption{Klick-Vergleich für Benutzererstellung}
    \label{tab:h6_klicks}
\end{table}

\subsection{Interpretation}

Ablauf sehr ähnlich, bei \ac{AWS} Cognito ist ein zusätzlicher Schritt zur Auswahl des User Pools erforderlich. Insgesamt sind 5 Klicks bei \ac{AWS} Cognito gegenüber 4 Klicks bei \ac{SAP IAS} notwendig.
$H_6$ wird nur teilweise bestätigt, da der UserPool auch direkt über einen Link geöffnet werden kann, wodurch sich die Klickanzahl auf 4 reduziert.
Zusätzlich bietet \ac{AWS} Cognito eine \ac{SDK} basierte Benutzererstellung, die den Klick-Aufwand weiter reduziert (selbst implementierbar).
Die Klicks beim Formular ausfüllen werden nicht gezählt, da diese lediglich Eingabefelder für E-Mail, Name usw. sind.


\section{H7: Recovery-Mechanismen}

Hypothese $H_7$ vergleicht die Sicherheit der Passwort-Reset-Mechanismen beider
\ac{IdP}s. Dies ist ein direkter \ac{IdP}-Vergleich.

\subsection{Messmethodik}

Der Passwort-Vergessen-Flow wird qualitativ und quantitativ verglichen:
\begin{itemize}
    \item \textbf{Mechanismus-Typen:} Link-basiert, Code-basiert, usw.
    \item \textbf{Sicherheitsaspekte:} Gültigkeitsdauer, Transportweg
    \item \textbf{Konfigurationsaufwand:} Schritte zur Aktivierung
\end{itemize}

\subsection{Messergebnisse}

\begin{table}[h]
    \centering
    \begin{tabular}{|l|c|c|}
        \hline
        \textbf{Aspekt} & \textbf{AWS Cognito} & \textbf{SAP IAS} \\
        \hline
        Mechanismus & 6-stelliger Code & Zeitlich limitierter Link \\
        \hline
        Gültigkeitsdauer & Standard: 1 Stunde & 1 Stunde bis 30 Tage \\
        \hline
        Zustellungsoptionen & E-Mail, SMS (kombinierbar) & E-Mail (Standard) \\
        \hline
        Alternative Methoden & Admin-Reset & Security Questions, PIN, Admin-Reset \\
        \hline
    \end{tabular}
    \caption{Vergleich der Recovery-Mechanismen}
    \label{tab:h7_recovery}
\end{table}

\ac{SAP IAS} bietet einen Forgot-Password-Flow mit konfigurierbarer Gültigkeitsdauer von 1 Stunde bis 30 Tagen für den Reset-Link. Der Standard-Mechanismus ist ein per E-Mail versendeter Link, alternativ können Security Questions (case-sensitive) oder ein PIN-Code konfiguriert werden, wobei SAP den PIN-Code nicht empfiehlt und der Nutzer diesen vorher setzen muss. Zusätzlich ist ein Admin-Reset mit temporärem Passwort möglich.

\ac{AWS} Cognito verwendet standardmäßig einen 6-stelligen Verifizierungscode, der per E-Mail oder SMS zugestellt wird. E-Mail und SMS können kombiniert werden, beispielsweise E-Mail als primärer Kanal mit SMS als Fallback. Die Gültigkeitsdauer beträgt standardmäßig 1 Stunde. Ein Admin-Reset mit temporärem Passwort ist ebenfalls verfügbar.

\subsection{Interpretation}

Bei \ac{SAP IAS} erfolgt der Passwort-Reset über einen zeitlich limitierten
Link, der per E-Mail versendet wird. Der Benutzer klickt auf den Link und wird
direkt zur Passwort-Änderungsseite weitergeleitet. Bei \ac{AWS} Cognito erhält
der Benutzer einen 6-stelligen Code per E-Mail, der manuell in die Anwendung
eingegeben werden muss. Link-basierte Resets gelten als sicherer, da
der Code nicht manuell übertragen werden muss und somit weniger anfällig für
Phishing oder Tippfehler ist. TODO belegen mit passender Literatur und citep


\section{H8: Integrierte Sicherheitsfeatures}

Hypothese $H_8$ quantifiziert die Anzahl standardmäßig integrierter
Sicherheitsfeatures.

\subsection{Feature-Checkliste}

\subsection{Feature-Checkliste}

\begin{table}[h]
    \centering
    \small
    \begin{tabular}{|l|c|c|}
        \hline
        \textbf{Feature} & \textbf{AWS Cognito}                           & \textbf{SAP IAS}                         \\
        \hline
        Rate-Limiting  & $\circ$ Konfigurierbar                           & $\bullet$ Standard                       \\
        \hline
        Brute-Force-Schutz & $\circ$ Konfigurierbar                           & $\bullet$ Standard                       \\
        \hline
        Audit-Logging  & $\bullet$ CloudWatch/CloudTrail (AWS Services)          & $\bullet$ Standard \\
        \hline
        Token-Expiration & $\bullet$ Konfigurierbar                           & $\bullet$ Konfigurierbar                 \\
        \hline
        E-Mail-Verifizierung & $\bullet$ Konfigurierbar                           & $\bullet$ Konfigurierbar                 \\
        \hline
        Passwort-Policies  & $\bullet$ Konfigurierbar                           & $\bullet$ Konfigurierbar                 \\
        \hline
        Account-Lockout  & $\circ$ Standard (5 Versuche)                            & $\bullet$ Standard                       \\
        \hline
        MFA-Unterstützung  & $\bullet$ Standard                           & $\bullet$ Standard                       \\
        \hline
        Risk-Based Authentication  & $\circ$ Über WAF                           & $\bullet$ Standard                       \\
        \hline
        Session-Timeout  & $\bullet$ Konfigurierbar                           & $\bullet$ Konfigurierbar (Standard: 12h) \\
        \hline
        \hline
        \textbf{Standard aktiviert}  & 2                            & 6                                        \\
        \hline
        \textbf{Konfigurierbar}  & 5                            & 4                                        \\
        \hline
        \textbf{Manuell/Extern}  & 1                            & 0                                        \\
        \hline
    \end{tabular}
    \caption{Feature-Checkliste für H8 - Aktualisiert}
    \label{tab:h8_features}
\end{table}

\textbf{Legende:} $\bullet$ = Standard/Verfügbar, $\circ$ = Konfiguration erforderlich

\subsection{Interpretation}

\ac{SAP IAS} bietet mehr Features standardmäßig aktiviert, während \ac{AWS} Cognito mehr Flexibilität durch Konfiguration ermöglicht.
Die Analyse zeigt signifikante Unterschiede in der Anzahl standardmäßig aktivierter Sicherheitsfeatures. \ac{SAP IAS} bietet mit 6 standardmäßig aktivierten Features gegenüber 2 bei \ac{AWS} Cognito eine deutlich höhere Out-of-the-Box-Sicherheit. Insbesondere sind bei \ac{SAP IAS} Rate-Limiting, Brute-Force-Schutz und Risk-Based Authentication ohne zusätzliche Konfiguration verfügbar.

\ac{AWS} Cognito erfordert für gleichwertige Sicherheitsfunktionen die Integration externer AWS-Services: Rate-Limiting und erweiterte Bedrohungserkennung werden über AWS WAF (Web Application Firewall) realisiert, Audit-Logging erfolgt über CloudWatch und CloudTrail. Dies erhöht die Konfigurationskomplexität und potenzielle Fehlerquellen.

Positiv hervorzuheben ist, dass \ac{AWS} Cognito in vielen Bereichen eine granularere Kontrolle durch Konfigurierbarkeit ermöglicht. Entwickler können die Sicherheitsparameter präzise an ihre Anforderungen anpassen, müssen dafür jedoch mehr Implementierungsaufwand investieren.

$H_8$ wird bestätigt: \ac{SAP IAS} bietet mit 6 standardmäßig aktivierten Sicherheitsfeatures eine um 200\% höhere Out-of-the-Box-Sicherheit als \ac{AWS} Cognito (2 Features). Für Unternehmen, die eine schnelle und sichere Implementierung ohne umfangreiche Konfiguration bevorzugen, ist \ac{SAP IAS} daher vorteilhaft.


\section{H9: Benutzerlöschung}

Hypothese $H_9$ vergleicht den Aufwand für die Löschung eines Benutzers über
die Admin-Konsolen beider \ac{IdP}s. Dies ist ein direkter \ac{IdP}-Vergleich.

\subsection{Messmethodik}

Die Benutzerlöschung wird in beiden Admin-Konsolen durchgeführt und die Anzahl
der Klicks/Schritte sowie die Anzahl der Bestätigungsdialoge gezählt.

\subsection{Messergebnisse}

\begin{table}[h]
    \centering
    \begin{tabular}{|l|c|c|}
        \hline
        \textbf{Schritt} & \textbf{AWS Cognito}          & \textbf{SAP IAS}          \\
        \hline
        1. Startseite der IdP öffnen & 1 & 1 \\
        \hline
        2. Users-Liste navigieren    & 1     & 1     \\
        \hline
        3. Benutzer suchen/auswählen & 1     & 1     \\
        \hline
        4. Löschen klicken   & 1     & 1     \\
        \hline
        5. Bestätigungsdialog  & 2     & 1     \\
        \hline
        \textbf{Gesamt Klicks} & 6     & 5     \\
        \hline
    \end{tabular}
    \caption{Klick-Vergleich für Benutzerlöschung}
    \label{tab:h9_klicks}
\end{table}

Bei \ac{AWS} Cognito sind 6 Klicks und 2 Bestätigungsdialoge erforderlich, während \ac{SAP IAS} 5 Klicks und 1 Bestätigungsdialog benötigt.
Ein wichtiger Unterschied ist, dass Beim Bestätigen muss zusätzlich der Nutzerzugriff deaktiviert werden, bevor der Nutzer gelöscht werden kann. Dadurch werden 2 Bestätigungsdialog-Klicks benötigt.
Bei \ac{AWS} Cognito wird zusätzlich als Ausgangspunkt der bereits ausgewählte Usepool genommen.


\section{Zusammenfassung und Diskussion}

\subsection{Bestätigte Hypothesen}

Die Evaluation zeigt, dass die Mehrheit der aufgestellten Hypothesen bestätigt werden konnte:

\begin{itemize}
    \item \textbf{$H_1$:} Bestätigt. Der deklarative Ansatz von \ac{SAP CAP} eliminiert die zyklomatische Komplexität der Autorisierungslogik vollständig (CC=26 $\rightarrow$ CC=0).

    \item \textbf{$H_2$:} Bestätigt. Das Hinzufügen einer neuen Rolle erfordert bei \ac{SAP CAP} Änderungen in nur 1 Datei (3 LOC), während bei Spring Boot 3-4 Dateien und 5 Methoden angepasst werden müssen.

    \item \textbf{$H_{3a}$:} Bestätigt. \ac{SAP CAP} reduziert die Token-Validierungs-Komplexität um 93\% (LOC) bzw. 100\% (CC), da das Framework die Validierung automatisch übernimmt.

    \item \textbf{$H_{3b}$:} Nicht bestätigt. Die Anbindung eines externen \ac{IdP} erfordert bei \ac{AWS} Cognito 6 Schritte, bei \ac{SAP IAS} 8 Schritte. \ac{AWS} bietet hier eine geringfügig einfachere Konfiguration.

    \item \textbf{$H_4$:} Bestätigt. \ac{SAP CAP} reduziert die Vendor-Kopplung um 60\% (2 Dateien vs. 5 Dateien mit Framework-spezifischen Imports).

    \item \textbf{$H_5$:} Bestätigt. \ac{SAP CAP} bietet vollautomatische Token-Validierung mit CC=0, während Spring Boot eine manuelle Implementierung mit CC=12 erfordert.

    \item \textbf{$H_6$:} Teilweise bestätigt. Die Benutzererstellung erfordert bei \ac{SAP IAS} 4 Klicks, bei \ac{AWS} Cognito 5 Klicks. Der Unterschied ist marginal.

    \item \textbf{$H_7$:} Bestätigt. \ac{SAP IAS} bietet mit dem Link-basierten Passwort-Reset einen sichereren Mechanismus als der Code-basierte Ansatz von \ac{AWS} Cognito.

    \item \textbf{$H_8$:} Bestätigt. \ac{SAP IAS} bietet 6 standardmäßig aktivierte Sicherheitsfeatures gegenüber 2 bei \ac{AWS} Cognito, was einer Steigerung von 200\% entspricht.

    \item \textbf{$H_9$:} Teilweise bestätigt. Die Benutzerlöschung erfordert bei \ac{SAP IAS} 5 Klicks, bei \ac{AWS} Cognito 6 Klicks. Der Unterschied ist marginal, wobei \ac{AWS} Cognito einen zusätzlichen Bestätigungsschritt erfordert.
\end{itemize}

\subsection{Gesamtvergleich}

\begin{table}[h]
    \centering
    \begin{tabular}{|l|c|c|}
        \hline
        \textbf{Kategorie} & \textbf{AWS Cognito + Spring Boot} & \textbf{SAP CAP + SAP IAS} \\
        \hline
        Sicherheits-LOC gesamt & $\sim$150 & $\sim$10 \\
        \hline
        Zyklomatische Komplexität (CC) & 26 & 0 \\
        \hline
        Dateien mit Security-Code  & 5 & 2 \\
        \hline
        Automatisierungsgrad & Manuell & Automatisch \\
        \hline
        Standard-Sicherheitsfeatures & 2 & 6 \\
        \hline
    \end{tabular}
    \caption{Gesamtvergleich der Implementierungen}
    \label{tab:gesamtvergleich}
\end{table}

Die Gesamtbetrachtung zeigt, dass \ac{SAP CAP} mit \ac{SAP IAS} in den meisten Framework-bezogenen Kategorien ($H_1$ bis $H_5$) deutliche Vorteile gegenüber Spring Boot mit \ac{AWS} Cognito bietet. Der deklarative Ansatz eliminiert die zyklomatische Komplexität vollständig und reduziert die Anzahl der sicherheitsrelevanten Dateien um 60\%. Die automatische Token-Validierung durch das Framework verhindert potenzielle Implementierungsfehler.

Im direkten \ac{IdP}-Vergleich ($H_6$ bis $H_9$) zeigt \ac{SAP IAS} Stärken bei den integrierten Sicherheitsfeatures und dem Recovery-Mechanismus. Die Unterschiede bei Benutzererstellung und -löschung sind marginal. Einzig bei der Unternehmensintegration ($H_{3b}$) bietet \ac{AWS} Cognito eine geringfügig einfachere Konfiguration.

\subsection{Limitationen der Studie}

Die Evaluation unterliegt folgenden Einschränkungen:

\begin{itemize}
    \item \textbf{Cloud-only bei SAP IAS:} Da \ac{SAP IAS} ausschließlich als Cloud-Service verfügbar ist, ist das \ac{SAP CAP} Backend auf die \ac{SAP BTP} deployt. Die Hypothesen ($H_6$ bis $H_9$) wurden daher in der Cloud-Umgebung / \ac{IdP}-UI gemessen.

    \item \textbf{Framework-Abhängigkeit:} Die Ergebnisse der Framework-Integration ($H_1$ bis $H_5$) sind spezifisch für die verwendeten Frameworks (Spring Boot 3.x, SAP CAP Java) und können bei anderen Versionen abweichen.

    \item \textbf{IdP-Vergleich:} Die Hypothesen $H_6$ bis $H_9$ messen nur die Admin-Konsolen-Erfahrung, nicht die programmatische API-Integration.
\end{itemize}
